% Автоматически перенесено из ПРИЛОЖЕНИЕ_2_ВЕРСИЯ_3.docx

% В оглавлении приложений отображаются только сами приложения; пункты 2.1,
% 3.1 и более глубокие уровни туда не выносятся. Команда находится внутри
% include, чтобы в .toc она была записана до заголовка приложения 2.
\addtocontents{toc}{\protect\setcounter{tocdepth}{0}}
\chapter{Словарь понятий и терминов}

В словарь включены исследовательские понятия и термины источников, необходимые для понимания проблематики и методики диссертации. Определения даны в том значении, в котором соответствующие понятия и обозначения используются в работе. Общеупотребительные исторические понятия и частные лексические соответствия, не требующие специального пояснения, в приложение не включены.

\section{I. Основные понятия правовой культуры и истории права}

\noindent\textbf{Вестготская правовая традиция} --- правовое наследие «Вестготской правды» (Lex Visigothorum), сохранявшее значение в каталонской книжной, судебной и документальной практике. Его присутствие в источнике не тождественно прямой текстуальной зависимости от конкретной рукописи «Вестготской правды».
\par\medskip

\noindent\textbf{Обычное право} --- совокупность признаваемых и применяемых в сообществе правовых обычаев и основанных на них способов регулирования отношений. Обычное право не тождественно неписаному: обычные нормы могли сохранять своё значение после письменной фиксации.
\par\medskip

\noindent\textbf{Письменная фиксация обычая} --- запись бытовавших обычаев, норм и процедур, предполагавшая их отбор, словесное оформление, редактирование и включение в структуру письменного источника. Такая фиксация рассматривается как переработка материала, а не как нейтральное воспроизведение ранее существовавшей нормы.
\par\medskip

\noindent\textbf{Правовая культура} --- исторически конкретная среда регулирования общественных отношений, в которой существовали и взаимодействовали нормы, представления о праве, способы правового действия и практики их применения. В диссертации это понятие объединяет обычай, язык правовых источников, судебную процедуру, положение участников и способы доказательства.
\par\medskip

\noindent\textbf{Правовая практика} --- совокупность конкретных действий и процедур, в ходе которых нормы и обычаи применялись, толковались, подтверждались или оспаривались. В работе она реконструируется прежде всего по актовому материалу и описаниям процедур в нормативных памятниках.
\par\medskip

\noindent\textbf{Правовая традиция} --- исторически воспроизводимая преемственность норм, представлений, терминов, формул и практик. Её передача могла сопровождаться отбором, переработкой и приспособлением унаследованного материала к новым условиям.
\par\medskip

\noindent\textbf{Правовое пространство} --- многомерная историческая среда существования права, включающая писаные и неписаные нормы, правовые представления, правовую культуру и правосознание, а также связи, возникающие в правовой деятельности людей и власти. Правовое пространство не сводится ни к правовой системе, ни к совокупности юрисдикций.
\par\medskip

\noindent\textbf{Правовой обычай} --- сложившееся в практике и признаваемое в конкретном сообществе правило поведения, которому придавалось нормативное значение. Обычай мог существовать без письменной записи и сохранять обычное происхождение после фиксации в тексте.
\par\medskip

\noindent\textbf{Римско-каноническая традиция} --- совокупность норм, понятий и способов правового рассуждения, восходивших к римскому и каноническому праву и распространявшихся через книжную традицию, обучение и деятельность профессионально подготовленных легистов.
\par\medskip

\noindent\textbf{Свод обычного права} --- письменный исторический источник, содержащий отобранную и упорядоченную запись местных обычаев, норм, процедур и правовых формул. Его текст является результатом отбора, формулирования, редактирования и систематизации материала.
\par\medskip

\noindent\textbf{Учёное право} --- традиция профессионального изучения, преподавания и толкования прежде всего римского и канонического права, связанная с университетской подготовкой и деятельностью легистов.
\par\medskip

\noindent\textbf{Юрисдикция} --- сфера властных, прежде всего судебных, полномочий конкретного носителя власти, определяемая кругом подвластных лиц, территорией и характером рассматриваемых дел.
\par\medskip

\section{II. Термины вычислительного анализа текстов}

\noindent\textbf{Граф текстуальных связей} --- модель, в которой тексты или их сегменты представлены как узлы, а выявленные между ними связи --- как рёбра; используется для сопоставления распределения текстовых сходств между памятниками.
\par\medskip

\noindent\textbf{Косинусная близость} --- мера сходства векторных представлений двух сегментов, отражающая степень совпадения их взвешенного лексического состава.
\par\medskip

\noindent\textbf{Локальное выравнивание Смита---Уотермана} --- метод поиска наиболее близкого локального фрагмента внутри двух последовательностей токенов. В работе применяется для выявления устойчивых оборотов и формул, занимающих лишь часть сравниваемых сегментов.
\par\medskip

\noindent\textbf{Мера Tesserae} --- мера текстовой близости, учитывающая совпадение редких и содержательно значимых токенов: чем реже общая единица встречается в корпусе, тем больший вес получает её совпадение.
\par\medskip

\noindent\textbf{Мягкая косинусная близость} --- мера текстовой близости, учитывающая наряду с точными совпадениями частичное формальное сходство словоформ, что позволяет работать с орфографической и морфологической вариативностью средневековых текстов.
\par\medskip

\noindent\textbf{Нормализация текста} --- приведение графических вариантов к сопоставимому аналитическому виду без исправления чтения источника; включает унификацию отдельных графем, удаление технического аппарата и стандартизацию записи.
\par\medskip

\noindent\textbf{Парето-слой} --- группа результатов, выделяемая на определённом этапе парето-фильтрации. Первый слой составляют пары, для которых нет другой пары, одновременно превосходящей их по всем учитываемым показателям; последующие слои определяются после удаления предыдущих.
\par\medskip

\noindent\textbf{Парето-фильтрация} --- отбор результатов по совокупности рассчитанных мер текстовой близости, при котором сохраняются пары, не уступающие другим одновременно по всем учитываемым показателям. Такой отбор позволяет не сводить оценку к одной заранее выбранной мере.
\par\medskip

\noindent\textbf{Перекрывающееся окно} --- фрагмент, полученный при дроблении крупного сегмента так, чтобы соседние фрагменты частично совпадали. Такой способ разбиения уменьшает риск потерять соответствие, проходящее через границу двух фрагментов.
\par\medskip

\noindent\textbf{Ранжирование} --- упорядочение отобранных пар с учётом их положения по нескольким мерам текстовой близости; определяет последовательность последующей содержательной проверки.
\par\medskip

\noindent\textbf{Редукция словоформ / стемминг} --- уменьшение вариативности словоформ путём осторожного отсечения отдельных окончаний и приведения форм к устойчивой основе. В работе используется вместо полной автоматической лемматизации.
\par\medskip

\noindent\textbf{Сегмент} --- структурно или технически выделенный фрагмент источника, выступающий единицей сопоставления: статья, глава, обычай, грамота либо часть более крупного текста.
\par\medskip

\noindent\textbf{Текстуальная связь} --- сходство между двумя фрагментами на уровне лексики, формулы, синтаксиса или структуры, которое может быть связано с заимствованием, общим источником или общей текстовой традицией. Вычислительный показатель сам по себе не доказывает прямого заимствования.
\par\medskip

\noindent\textbf{Токенизация} --- разбиение текста на минимальные формально выделяемые единицы --- токены, с которыми затем выполняются нормализация, фильтрация и сопоставление.
\par\medskip

\noindent\textbf{TF--IDF} --- способ взвешивания токенов по сочетанию их частоты в данном сегменте и редкости во всём корпусе. В работе используется на этапе предварительного отбора возможных соответствий.
\par\medskip

\section{III. Язык права, речевые формулы и перевод}

\noindent\textbf{Вернакуляр} --- местный язык, противопоставленный в данном контексте книжной латыни. Для основной части исследуемого материала таким вернакуляром является старокаталанский язык.
\par\medskip

\noindent\textbf{Вернакуляризация} --- расширение употребления местного романского языка в сферах, где прежде преобладала латынь, в том числе в правовой письменности, переводах и записи местного обычая.
\par\medskip

\noindent\textbf{Генерализация при переводе} --- лексико-семантическая трансформация, при которой единица исходного текста передаётся словом или выражением с более широким значением.
\par\medskip

\noindent\textbf{Глаголы речи} --- глаголы, обозначающие речевые действия участников правовой процедуры: заявление, ответ, обвинение, наименование, клятву и другие формы высказывания.
\par\medskip

\noindent\textbf{Дефиниция в правовом тексте} --- содержащееся в самом источнике разъяснение значения слова, термина или понятия, нередко оформленное специальной пояснительной конструкцией. Такая дефиниция отражает понятийную работу составителя текста и не тождественна современному исследовательскому определению.
\par\medskip

\noindent\textbf{Диглоссия} --- социально закреплённое распределение функций между языковыми кодами, различавшимися сферами употребления и престижем. Для исследуемого материала понятие применяется к функциональному соотношению латыни и старокаталанского языка и не тождественно индивидуальному билингвизму.
\par\medskip

\noindent\textbf{Заимствование} --- переход языковой единицы из одного языка в другой и её включение в новую языковую среду, с адаптацией или без неё. В работе рассматриваются прежде всего романские элементы в латинских текстах и латинские элементы в старокаталанских.
\par\medskip

\noindent\textbf{Иноязычное включение} --- слово, словосочетание или более протяжённый фрагмент на языке, отличном от основного языка текста. Степень его графической, морфологической и синтаксической интеграции может быть различной.
\par\medskip

\noindent\textbf{Конкретизация при переводе} --- лексико-семантическая трансформация, при которой единица исходного текста передаётся словом или выражением с более узким и конкретным значением.
\par\medskip

\noindent\textbf{Латинизм} --- латинская языковая единица или восходящая к латыни форма, сохраняемая в старокаталанском тексте и связанная с книжной или латинской правовой традицией.
\par\medskip

\noindent\textbf{Модуляция при переводе} --- переводческая трансформация, при которой исходное содержание передаётся через логически связанное с ним выражение, представляющее ту же ситуацию с иной смысловой точки зрения.
\par\medskip

\noindent\textbf{Мультилингвизм} --- сосуществование и взаимодействие нескольких языков или языковых кодов в одной социальной и письменной среде. Понятие шире диглоссии и не предполагает обязательного устойчивого распределения функций между всеми используемыми языками.
\par\medskip

\noindent\textbf{Объяснительный перевод} --- передача единицы исходного текста более развёрнутым выражением, раскрывающим её значение вместо краткого прямого соответствия.
\par\medskip

\noindent\textbf{Перформатив} --- высказывание, произнесение которого при соблюдении установленных условий само составляет правовое действие или необходимую часть такого действия.
\par\medskip

\noindent\textbf{Прямая речь в правовом тексте} --- фрагмент, оформленный как реплика участника либо как образец слов, которые надлежало произнести в определённой правовой ситуации. Такая запись не обязательно воспроизводит реально произнесённые слова дословно.
\par\medskip

\noindent\textbf{Речевая формула} --- устойчивое выражение, связанное с определённым правовым действием или этапом процедуры и воспроизводимое в источнике как образец требуемого или ожидаемого высказывания.
\par\medskip

\noindent\textbf{Социальная маркировка словоупотребления} --- прямое или косвенное указание в источнике на связь определённого слова или выражения с конкретной социальной группой или речевой средой.
\par\medskip

\noindent\textbf{Терминологическое соответствие} --- соотношение обозначений одного или близких понятий в разных языковых версиях текста. Такое соответствие может быть устойчивым, частичным или контекстуальным и не предполагает полного совпадения значений.
\par\medskip

\noindent\textbf{Формуляр} --- устойчивая структура документа и последовательность повторяющихся словесных формул, посредством которых оформлялось определённое правовое действие.
\par\medskip

\noindent\textbf{Цитирование} --- воспроизведение в правовом тексте узнаваемого фрагмента другого, обычно авторитетного, текста. В отличие от более свободного заимствования цитирование сохраняет текстовую узнаваемость источника.
\par\medskip

\noindent\textbf{Язык права} --- лексика, синтаксические конструкции и устойчивые формулы, используемые для выражения правовых понятий, норм, действий и процедур.
\par\medskip

\section{IV. Судебная процедура, доказательство и документ}

\noindent\textbf{«Божий суд» / ордалия (лат} --- iudicium Dei) --- судебное испытание, исход которого истолковывался как обнаружение Божьей воли относительно правоты участника спора или обвинения. К ордалиям относились испытания водой и железом; двусторонней формой испытания выступал судебный поединок.
\par\medskip

\noindent\textbf{Клятва «по иудейскому обычаю» (лат} --- more iudaico) --- предписанный иудею порядок принесения судебной клятвы, учитывавший конфессиональную принадлежность клянущегося и отличавшийся от христианской формы.
\par\medskip

\noindent\textbf{Окончательный отказ от притязаний (лат} --- diffinitio / definitio) --- акт окончательного прекращения требований по спорному праву, закреплявший завершение конфликта и отказ от последующих притязаний по тому же основанию.
\par\medskip

\noindent\textbf{Перечень жалоб и притязаний (лат} --- querimonia, мн. querimoniae) --- запись жалоб и требований стороны, возникших в связи с конфликтом. Такой текст фиксировал позицию участника и мог предшествовать судебному разбирательству, арбитражу или переговорам.
\par\medskip

\noindent\textbf{Письменный акт (лат} --- instrumentum) --- документ, оформлявший или удостоверявший правовое действие, обязательство либо распоряжение. В источниках термин мог уточняться указанием на публичный характер или предмет акта.
\par\medskip

\noindent\textbf{Пригодный свидетель (лат} --- idoneus testis) --- свидетель, признававшийся допустимым и заслуживающим доверия с учётом требований к его положению, репутации, отношениям со сторонами и обстоятельствам дела.
\par\medskip

\noindent\textbf{Признание (лат} --- recognitio) --- акт признания права или требования другой стороны, нередко составлявшийся как результат судебного либо договорного разрешения конфликта.
\par\medskip

\noindent\textbf{Примирение (лат} --- pacificatio) --- акт или формула, закреплявшие восстановление мира между сторонами после конфликта. В актовом материале pacificatio могла сочетаться с признанием права и отказом от притязаний.
\par\medskip

\noindent\textbf{Публичная грамота (старокат} --- carta pública) --- письменный акт, доказательное значение которого связывалось с установленным порядком составления и удостоверения. В «Обычаях Тортосы» публичная грамота названа наряду со свидетелями и признанием стороны среди способов подтверждения обстоятельств.
\par\medskip

\noindent\textbf{Публичное удостоверение} --- придание документу публично признаваемой достоверности через установленный порядок его составления и засвидетельствования писцом, нотариусом, судом или канцелярией, а также через подписи, знаки и иные формы удостоверения.
\par\medskip

\noindent\textbf{Публичный акт о долге (лат} --- publicum crediti instrumentum) --- удостоверенный письменный акт, фиксировавший долговое обязательство и его содержание.
\par\medskip

\noindent\textbf{Репутация / «слава» (лат} --- fama; старокат. --- bona fama, mala fama, d’íntegra fama) --- сложившееся в сообществе мнение о человеке, которое могло учитываться при оценке его слова и процессуального положения. Источники различают добрую, дурную и безупречную славу.
\par\medskip

\noindent\textbf{Свидетельское показание (лат} --- testimonium) --- сообщение свидетеля об обстоятельствах дела, приобретавшее доказательное значение при соблюдении требований к самому свидетелю и процедуре его опроса.
\par\medskip

\noindent\textbf{Судебная клятва (лат} --- sacramentum, iurare / jurare; старокат. --- sagrament, jurar) --- ритуализованное словесное действие, имевшее доказательное или процессуальное значение. Существительные sacramentum и sagrament обозначали клятву как акт или обязательство, глаголы iurare / jurare и jurar --- её совершение; конструкции fer, donar, pendre и tornar sagrament фиксировали различные действия с клятвой в судебной процедуре.
\par\medskip

\noindent\textbf{Судебное доказательство} --- признаваемый процедурой способ подтверждения или опровержения обстоятельств, значимых для разрешения дела. В исследуемых источниках к доказательственным практикам относились клятва, свидетельские показания, ордалия, судебный поединок, репутационные критерии и письменные акты.
\par\medskip

\noindent\textbf{Судебный поединок (старокат} --- batalia) --- двустороннее судебное испытание, при котором стороны или их представители вступали в ритуализованное вооружённое состязание, а его исход имел доказательное значение.
\par\medskip

\noindent\textbf{Уступка / освобождение от притязаний (лат} --- evacuatio) --- акт отказа стороны от спорного права или требования в пользу другой стороны. Evacuatio могла оформлять один из завершающих этапов судебного или договорного урегулирования конфликта.
\par\medskip

\section{V. Социальная структура и имущественные отношения}

\noindent\textbf{Арсия (лат} --- arcia) --- платёж или взыскание в пользу сеньора, связывавшееся с пожаром в крестьянском хозяйстве и входившее в комплекс «дурных обычаев».
\par\medskip

\noindent\textbf{Бакаларии (лат} --- baccalarii) --- в рассматриваемом каталонском контексте --- молодые люди крестьянского происхождения, включённые в свиту сеньора и участвовавшие в военных действиях, но не имевшие рыцарского статуса.
\par\medskip

\noindent\textbf{Брачное обеспечение / брачный дар (лат} --- sponsalicium; старокат. --- espoali) --- имущественное предоставление в пользу жены, связанное с заключением брака; в старокаталанской версии Recognoverunt proceres обозначается термином espoali.
\par\medskip

\noindent\textbf{Вальвассоры (лат} --- valvassores) --- слой каталонской знати, занимавший положение между верхушкой аристократии и рыцарством и включённый в иерархию держаний и службы.
\par\medskip

\noindent\textbf{Владение / держание (лат., старокат} --- honor) --- обозначение владения или держания, включавшего землю и связанные с ней права и доходы.
\par\medskip

\noindent\textbf{«Дурные обычаи» (старокат} --- mals usos) --- комплекс сеньориальных повинностей и взысканий, характеризовавших зависимость части крестьян Старой Каталонии; к ним относились remença personal, intestia, exorquia, cugucia, arcia и ferma de spoli.
\par\medskip

\noindent\textbf{Интестия (лат} --- intestia) --- сеньориальное притязание на часть имущества зависимого крестьянина, умершего без завещания.
\par\medskip

\noindent\textbf{Кастеляны (лат} --- castlani) --- держатели или управители замков, связанные с местной сеньориальной властью, военной организацией и осуществлением юрисдикционных полномочий.
\par\medskip

\noindent\textbf{Клятва верности / присяга (лат} --- sacramentum fidelitatis) --- клятвенное обещание верности, связанное с отношениями личной зависимости и оммажем. В отличие от судебной клятвы, такая присяга не выступала способом судебного доказательства.
\par\medskip

\noindent\textbf{Комиторы (лат} --- comitores) --- группа каталонской знати, занимавшая положение ниже верхушки магнатов и связанная с местной системой держаний и службы.
\par\medskip

\noindent\textbf{Кугусия (лат} --- cugucia / cugutia) --- сеньориальное имущественное притязание, возникавшее в связи с супружеской неверностью жены зависимого крестьянина и входившее в комплекс «дурных обычаев».
\par\medskip

\noindent\textbf{Леуда (старокат} --- lleuda; лат. --- leuda) --- торговый сбор, взимавшийся с товаров или торговых операций; право на его получение могло составлять часть доходов с владения или держания.
\par\medskip

\noindent\textbf{Линьяж (старокат} --- linatge) --- родственная группа, связывавшая своё происхождение с общим предком; в правовой и политической лексике XIII~в. термин прежде всего обозначал знатный род и его потомков.
\par\medskip

\noindent\textbf{Магнаты земли (лат} --- magnates terrae) --- обозначение верхушки каталонской знати, к которой относились наиболее влиятельные графы, виконты и крупные держатели замков, обладавшие значительными владениями и властными полномочиями.
\par\medskip

\noindent\textbf{Манс / крестьянское хозяйство (лат} --- mansus; старокат. --- mas) --- хозяйственно-земельная единица, включавшая дом, землю и связанные с ними права и повинности; могла выступать объектом держания, передачи и судебного спора.
\par\medskip

\noindent\textbf{Приданое (лат} --- dos; старокат. --- exovar) --- имущество, предоставлявшееся жене в связи с заключением брака; старокаталанский термин exovar соответствовал латинскому dos.
\par\medskip

\noindent\textbf{Приращение к приданому (старокат} --- creix) --- часть брачного имущественного обеспечения; в Recognoverunt proceres creix составлял треть совокупной суммы exovar и espoali.
\par\medskip

\noindent\textbf{Ременса (старокат} --- remença personal) --- выкуп, который зависимый крестьянин должен был уплатить при уходе с держания или смене сеньора; от названия этой повинности произошло обозначение ременсов.
\par\medskip

\noindent\textbf{Ременсы (старокат} --- remences) --- зависимые крестьяне Старой Каталонии, положение которых связывалось с личной зависимостью от сеньора и обязанностью выкупа при уходе с держания или смене сеньора.
\par\medskip

\noindent\textbf{Фирма д’эсполи форсада (ferma de spoli)} --- сеньориальный платёж, связанный с браком зависимого крестьянина или передачей хозяйства и входивший в комплекс «дурных обычаев».
\par\medskip

\noindent\textbf{Экзоркия (лат} --- exorquia) --- сеньориальное право на имущество зависимого крестьянина, умершего без наследников, входившее в комплекс «дурных обычаев».
\par\medskip
